%! AP Statistics — Unit 9, Lesson 9.1 — Group Activity (Answer Key)
\documentclass[10pt]{article}
\usepackage{apstats-article}
\usepackage{apstats-key}

\begin{document}

\pageheader{Unit 9, Lesson 9.1}{Group Activity --- Answer Key}
\namepartnerperiod

\vspace{0.1in}
\textbf{Is the Slope Real? Exploring Sampling Variability in Regression}

\vspace{0.05in}
Below are three regression scenarios. In each, a researcher collected data and fit a
least-squares regression line. Your job: decide whether the observed slope seems
convincingly different from 0, or whether it could plausibly be due to random variation.

\vspace{0.1in}
\begin{center}
\renewcommand{\arraystretch}{1.4}
\begin{tabular}{clccc}
    \textbf{\#} & \textbf{Context} & \textbf{Explanatory ($x$)} & \textbf{Response ($y$)} & \textbf{$n$, $b$} \\
    \hline
    1 & Sunlight and plant growth & hours of sunlight & growth (cm) & $n=200$,\ $b=0.45$ \\
    2 & Driving and gas use       & miles driven      & gallons used & $n=8$,\ $b=1.2$ \\
    3 & Shoe size and GPA         & shoe size         & GPA & $n=40$,\ $b=0.03$ \\
\end{tabular}
\end{center}

\vspace{0.15in}

\begin{tcolorbox}[colback=white, colframe=black!40,
    title=\textbf{Tier R --- Remediate (Key)},
    fonttitle=\bfseries, arc=2mm, left=3mm, right=3mm, top=2mm, bottom=2mm]

\textbf{Scenario 1 only:} $n = 200$,\ $b = 0.45$ (hours of sunlight vs.\ plant growth in cm)

\begin{enumerate}[label=\alph*.]

  \item \ans{For each additional hour of sunlight, the predicted plant growth increases by 0.45 cm.}

        \vspace{0.05in}

  \item \textbf{Very surprising} (circled). \ans{With $n = 200$ observations, a slope of
        $b = 0.45$ represents a consistent trend across many data points. If $\beta = 0$,
        it would be very unlikely to observe a slope this large just by chance.}

        \vspace{0.05in}

  \item \ans{If $\beta = 0$, that means hours of sunlight have no real linear effect on
        plant growth in the population --- knowing the sunlight hours does not help
        predict plant growth.}

\end{enumerate}
\end{tcolorbox}

\begin{tcolorbox}[colback=white, colframe=black!40,
    title=\textbf{Tier A --- Approaching Proficiency (Key)},
    fonttitle=\bfseries, arc=2mm, left=3mm, right=3mm, top=2mm, bottom=2mm]

\begin{enumerate}[label=\alph*.]

  \item Slope interpretations:

        Scenario 1: \ans{For each additional hour of sunlight, the predicted plant growth increases by 0.45 cm.}

        Scenario 2: \ans{For each additional mile driven, the predicted gas used increases by 1.2 gallons.}

        Scenario 3: \ans{For each one-unit increase in shoe size, the predicted GPA increases by 0.03.}

  \item Convincing evidence that $\beta \ne 0$?

        Scenario 1: \ans{Yes --- convincing. The slope $b = 0.45$ is not tiny, and the
        large sample size ($n = 200$) means sampling variability is low. A slope this far
        from 0 is unlikely by chance alone.}

        Scenario 2: \ans{Somewhat --- but not convincing. Although $b = 1.2$ seems large,
        the sample is only $n = 8$. A single unusual observation could drive the slope.
        High variability with small $n$ means this $b$ could easily occur by chance.}

        Scenario 3: \ans{Not convincing. The slope $b = 0.03$ is very small. Even with
        $n = 40$, a slope this close to 0 could easily arise from random variation when
        $\beta = 0$.}

  \item Most confident: \ans{Scenario 1} \quad Least confident: \ans{Scenario 3 (or 2)}

        Reasoning: \ans{Scenario 1 has a moderate slope \emph{and} a large sample --- both
        factors support convincing evidence. Scenario 3 has a very small slope, so it is
        hard to distinguish from $\beta = 0$ even at $n = 40$.}

\end{enumerate}
\end{tcolorbox}

\begin{tcolorbox}[colback=white, colframe=black!40,
    title=\textbf{Tier E --- Extension (Key)},
    fonttitle=\bfseries, arc=2mm, left=3mm, right=3mm, top=2mm, bottom=2mm]

\begin{enumerate}[label=\alph*., start=4]

  \item \ans{Both the slope magnitude \emph{and} the sample size matter. A large slope
        from a very small sample ($n = 8$) is less convincing than a moderate slope from
        a large sample ($n = 200$), because small samples have high sampling variability.
        An unusual slope from 8 points could easily be driven by one or two atypical
        observations. With 200 points, that same general trend is much harder to explain
        by chance alone.}

  \item \ans{Sampling variability always produces non-zero slopes, even when $\beta = 0$.
        If we drew many random samples from a population with no linear relationship
        ($\beta = 0$), the sample slopes $b$ would vary around 0 but would almost never
        equal 0 exactly. So observing $b \ne 0$ tells us nothing by itself --- we need to
        know how unlikely our specific $b$ would be if $\beta$ truly were 0.}

\end{enumerate}
\end{tcolorbox}

\end{document}
